\chapter{引言}
\label{ch:introduction}

\section{长上下文大语言模型的兴起}
大语言模型（LLMs）已经迅速成为研究与生产系统中的核心工作负载。以 LLaMA 和 Qwen 为代表的近期模型表明，模型能力的扩展不仅体现在参数规模上，也体现在有效上下文窗口的持续扩展上，使模型能够处理更长的文档、更大的代码仓库以及持续多轮的对话历史~\cite{llm_context, ring_attention}。过去只局限于几千 token 的上下文，如今已经常见于数万甚至数十万 token 的范围，且研究原型仍在继续向更长上下文推进。

这一变化也改变了推理阶段的主导瓶颈。在中等序列长度下，推理常被视为一个稠密线性代数问题；而在长上下文下，系统越来越表现为一种内存密集型流式工作负载，其中 KV 缓存的存储、搬移与同步成为一阶问题。

\section{生成式内存墙}
LLM 推理通常包含两个硬件行为不同的阶段：\textit{prefill} 与 \textit{decode}。Prefill 可以并行处理输入提示词，较好地映射到现代加速器的高吞吐矩阵引擎上；而 decode 则是串行的。每生成一个新 token，当前查询向量都需要与历史所有 token 对应的键和值进行注意力计算，这些历史状态保存在“键值缓存（KV cache）”中~\cite{vllm}。

随着序列长度增长，KV 缓存逐渐成为解码阶段跨步保留的主要状态。对于一个 7B 级别模型，在 $128{,}000$ token 上下文下，单个请求的 KV 缓存就可能占用数 GB 空间，尚未计入并发用户和批处理开销~\cite{vllm}。在实际服务系统中，这会很快超出片上高带宽内存（HBM）的容量，使长上下文推理从一个纯计算问题转变为同时受制于内存容量和内存带宽的问题。

为扩展容量，系统设计者越来越多地依赖片外内存扩展机制，包括基于 Compute Express Link（CXL）的内存层级~\cite{cxl_standard}。但仅仅扩展容量，并不能消除 decode 阶段反复搬运 KV 缓存所带来的带宽和延迟开销。

\section{现有基础设施中的挑战}

\subsection{CXL 与带宽瓶颈}
Compute Express Link 通过将远端内存设备映射进主机地址空间，使内存池化和分层内存扩展成为可能~\cite{cxl_pond}。这有助于缓解 KV 缓存的容量问题，但也暴露出新的吞吐瓶颈。在传统的主机聚合式设计中，主机或加速器在每一步 decode 时仍需跨互联链路拉取远端 KV 张量。当上下文较长时，这种对稠密张量的反复搬运会迅速占满链路带宽，并在 CXL 交换机端口形成排队，进一步抬高平均延迟与尾延迟~\cite{distserve}。

\subsection{存内计算与格式失配}
存内计算（PIM）是应对互联饱和的一个有吸引力的方向，因为它将计算移动到 KV 缓存数据附近~\cite{samsung_pim}。困难在于，现代 LLM 服务栈常常以 INT8 或 INT4 等量化格式保存 KV 缓存，以提高有效内存容量~\cite{llm_outliers, gptq}；而许多现有 PIM 设计却默认采用浮点非线性数据通路。二者若直接结合，就必须反复执行反量化、浮点 Softmax 计算和重新量化。由此带来的格式转换开销，会削弱近内存执行原本希望获得的效率优势。

\section{研究动机与目标}
本文的核心动机很直接：长上下文解码的主要限制因素，不再是原始算术吞吐，而是 KV 缓存能否被高效地存储、搬运和同步。如果注意力计算能够在 KV 数据附近完成，而分布式归一化又能通过轻量级归约而不是主机搬运大张量实现，那么单节点和多节点两类瓶颈都有望被同时缓解。

因此，本文有两个目标。第一，设计一种单节点近内存注意力引擎，在保持量化数据通路的同时避免浮点格式转换。第二，将这一思路扩展到分离式 CXL 内存池中，以支持可扩展且对延迟敏感的多租户服务。

\section{论文贡献}
本文的主要贡献分为单节点和分布式两部分：

\subsection{LKC-CXL-PIM（单节点体系结构）}
我们设计了一种面向单节点长上下文解码的 CXL 附加 PIM 体系结构。
\begin{itemize}
    \item \textbf{整数非线性单元（iNLU）：} 提出一种面向整数域的非线性数据通路，通过范围约简与紧凑多项式流水线逼近 Softmax 中的指数函数，从而避免在近内存路径中使用浮点 Softmax 模块。
    \item \textbf{异常值感知路由：} 设计双路径执行机制，使绝大多数激活经过高密度 INT8 阵列处理，而少量大幅值激活被重定向到小型高精度溢出路径。
    \item \textbf{单节点收益验证：} 基于周期级内存仿真与综合代理估计，证明该设计在 128K 上下文长度下可将读取延迟降低 98.9\% 以上，并相对于基线实现 17\% 的整体逻辑层面积缩减。
\end{itemize}

\subsection{DisaggKV（多节点服务系统）}
我们进一步将该架构扩展到面向多租户服务的分离式 CXL 内存池。
\begin{itemize}
    \item \textbf{局部性感知放置：} 提出一种区分共享提示页和私有 KV 页的放置策略，使高复用前缀能够被局部保留，减少不必要的数据移动。
    \item \textbf{基于 CXL 的点对点归约：} 以端点间交换紧凑的 \texttt{Global\_Max} 和 \texttt{Global\_Sum} 取代主机聚合，从而将总互联流量降低约 84\%，并将面向主机的流量降低 90\% 以上。
    \item \textbf{尾延迟与恢复能力提升：} 在分布式评估中，DisaggKV 可在大约 2.9K requests/s 之前维持 50\,ms 的尾延迟目标，并在故障注入模型下实现亚秒级恢复。
\end{itemize}

\section{论文结构}
本文其余内容安排如下。第~\ref{ch:background} 章回顾 LLM 推理、CXL 内存系统、排队行为和 PIM 的技术背景。第~\ref{ch:motivation} 章分析提出架构所针对的关键瓶颈。第~\ref{ch:inlu_architecture} 章介绍单节点的 LKC-CXL-PIM 设计，第~\ref{ch:disaggkv_system} 章进一步扩展为分布式 DisaggKV 系统。第~\ref{ch:methodology} 章说明仿真与分析方法，第~\ref{ch:evaluation} 章给出对应实验结果。第~\ref{ch:related_work} 章将本文置于相关研究背景下讨论，第~\ref{ch:conclusion} 章总结全文并展望未来工作。
